\section{Introducción}


Los transformadores monofásicos desempeñan un rol crítico en los sistemas eléctricos modernos, ya que permiten la transferencia eficiente de energía entre niveles de tensión distintos mediante el fenómeno de inducción electromagnética. Para garantizar una operación óptima y predecir con exactitud el comportamiento de estas máquinas bajo diferentes perfiles de carga, resulta esencial desarrollar un modelo de circuito equivalente que represente fielmente sus características eléctricas y los fenómenos magnéticos internos.

Esta práctica de laboratorio se centra en la caracterización integral de un transformador monofásico a través de una metodología experimental que incluye ensayos de resistencia de devanados, determinación de la relación de transformación, pruebas de polaridad, así como los ensayos de vacío y cortocircuito. El análisis técnico de los datos recabados en estas pruebas permitirá cuantificar los parámetros del modelo tales como resistencias, reactancias de dispersión y los componentes de la rama de excitación, facilitando el cálculo preciso de indicadores clave como la eficiencia y la regulación de tensión.

